% !PRINTED 392
\weberpage{392}
2. \gesperrt{Ist $p$ eine in $D$ aufgehende ungerade Primzahl, so ist $a$ dann und nur dann Normenrest, wenn $a$ quadratischer Rest von $p$ ist.}

Das ergibt sich unmittelbar aus der Bedingung für die Normenreste:

\weq{4}{x^2 - D y^2 \equiv 4a \pmod{p}},

die sich in diesem Fall auf $x^2 \equiv 4a \pmod{p}$ reduziert.

Die Primzahl 2 kommt nicht in Betracht, weil $x^2 \equiv a \pmod{2}$ immer lösbar ist, wohl aber die Moduln 4 und 8. Ist $D$ ungerade, so ist $D \equiv 1 \pmod{4}$, und die Bedingung für einen Normenrest $a$ ist:

\[ \frac{x^2 - D y^2}{4} \equiv a \pmod{4} \text{ oder } \pmod{8}, \]

und diese Kongruenz ist für jedes ungerade $a$ lösbar (durch gerade $x, y$), also:

3. \gesperrt{Geht 4 in $D$ nicht auf, so ist jede ungerade Zahl $a$ Normenrest von 4 und von 8.}

Ist $D$ durch 4 teilbar, so ist die Bedingung für einen Normenrest $a$ die Möglichkeit der Kongruenz:

\weq{5}{x^2 - \frac{D}{4} y^2 \equiv a \pmod{4} \text{ oder } \pmod{8}.}

Da $a$ ungerade ist, so können $x$ und $D y^2/4$ nicht beide ungerade sein.

Geht man hiernach die einzelnen Fälle durch, so ergibt sich

\[
\begin{array}{l}
\dfrac{D}{4} \equiv 0 \pmod{8}, \quad a \equiv 1 \pmod{4}, \quad a \equiv 1 \pmod{8},\\[6pt]
\dfrac{D}{4} \equiv 1 \pmod{8}, \quad \text{keine Bedingung für } a,\\[6pt]
\dfrac{D}{4} \equiv 2 \pmod{8}, \quad a \equiv +1,\ -1 \pmod{8},\\[6pt]
\dfrac{D}{4} \equiv 3 \pmod{8}, \quad a \equiv 1 \pmod{4}, \quad a \equiv +1,\ 5 \pmod{8},\\[6pt]
\dfrac{D}{4} \equiv 4 \pmod{8}, \quad a \equiv 1 \pmod{4}, \quad a \equiv +1,\ 5 \pmod{8},\\[6pt]
\dfrac{D}{4} \equiv 5 \pmod{8}, \quad \text{keine Bedingung für } a,\\[6pt]
\dfrac{D}{4} \equiv 6 \pmod{8}, \quad a \equiv +1,\ 3 \pmod{8},\\[6pt]
\dfrac{D}{4} \equiv 7 \pmod{8}, \quad a \equiv 1 \pmod{4}, \quad a \equiv +1,\ 5 \pmod{8}.
\end{array}
\]
